\documentclass[11pt]{article}
\usepackage[margin=25mm]{geometry}
\usepackage{fontspec}
\setmainfont{Latin Modern Roman}
\usepackage{amsmath,amssymb,mathtools}
\usepackage{unicode-math}
\setmathfont{Latin Modern Math}
\usepackage{xurl,hyperref}
\hypersetup{colorlinks=true,urlcolor=blue}
\setcounter{secnumdepth}{0}
\setlength{\parindent}{0pt}
\setlength{\parskip}{5pt}
\setlength{\emergencystretch}{3em}
\newcommand{\sref}[2]{\hyperlink{source:#1:#2}{[S#2]}}
\newcommand{\cataloguescope}[1]{\paragraph{Recorded target.}#1}
\newcommand{\R}{\mathbb R}
\newcommand{\C}{\mathbb C}
\newcommand{\N}{\mathbb N}
\newcommand{\Z}{\mathbb Z}
\newcommand{\Q}{\mathbb Q}
\newcommand{\F}{\mathbb F}
\newcommand{\eps}{\varepsilon}
\begin{document}
\section{Group-Valued Lévy Processes with Type-II Arveson Systems}\label{statement}

\subsection{Definitions and mathematical statement}

\paragraph{Shared notation.}$\mathbb N=\{1,2,\ldots\}$; the standard number fields have their usual meanings. All additional conventions are stated locally.


\paragraph{Definitions and conventions.} A group-valued Lévy process starts at the identity, has stationary independent increments interpreted through group multiplication, and is stochastically continuous in the chosen group topology. Increment probability spaces give a continuous product of Hilbert spaces. A unit is a multiplicative measurable section; the associated Arveson system is type I when generated by units and type II when it has units whose generated subsystem is proper. A homeomorphism group uses its specified topology, so the group and topology are part of the model.

\paragraph{Source scope and qualifications.} Freshly consulted primary publication evidence gives a positive answer to the broad group-valued existence question: Floricel--Wadel construct type-II\_infinity systems from infinite-dimensional Banach-space-valued pure-jump Lévy processes, and an additive Banach space is a topological group. This does not answer the narrower suggested homeomorphism-group example. The saved partially\_solved status is retained as provenance rather than changed.

\paragraph{Catalogue mathematical target.}
\paragraph{Catalogue wording retained for traceability.}
Do there exist Lévy processes with values in a group, for example a group of homeomorphisms of the circle, whose Arveson system is of type $\mathrm{II}$ rather than type $\mathrm I$?

\paragraph{Saved literature context.} The supplied catalogue's assessment dated 2026-08-21 records: A 2025 primary paper constructs nonisomorphic type II product systems from pure-jump Levy processes on an infinite-dimensional Banach space, which is an additive topological group; the narrower homeomorphism-group example in the extracted question is not resolved. The assessment is retained as dated provenance; it is not a new comprehensive literature or open-status certification.
\paragraph{Scope and source qualification.}Freshly consulted primary publication evidence gives a positive answer to the broad group-valued existence question: Floricel--Wadel construct type-II_infinity systems from infinite-dimensional Banach-space-valued pure-jump Lévy processes, and an additive Banach space is a topological group. This does not answer the narrower suggested homeomorphism-group example. The saved partially_solved status is retained as provenance rather than changed.
\subsection{Short English statement}
Find a group-valued Lévy process with a type-II product system, including the suggested case of circle homeomorphisms.
\subsection{Research attempt}
\paragraph{Model, date and investigated scope.}
OpenAI GPT-6 Astra (configured model \texttt{gpt-6-astra}), ultra reasoning; generated 11 October 2026. Proves that the elementary integer-valued Poisson candidate is unit-generated, while retaining the source-reported broad positive result.

\paragraph{Sources and literature check.}
The complete supplied definition was reviewed, including its Floricel–Wadel 2025 positive-result qualification and the distinction from circle homeomorphisms. This attempt did not independently consult or reprove that paper. The elementary Poisson calculation below claims no novelty and does not alter the saved source status.

\paragraph{Context and investigated scope.} The complete supplied definition was reviewed, including its Floricel–Wadel 2025 positive-result qualification and the distinction from circle homeomorphisms. This attempt did not independently consult or reprove that paper. The elementary Poisson calculation below claims no novelty and does not alter the saved source status.

\paragraph{A finite-rate Poisson candidate gives type I.} Let $N_t$ be a rate-$\lambda>0$ Poisson process valued in the additive group $\mathbb Z$, and take the Hilbert space of its increment sigma-field on $(0,t)$. For every real $\theta$, the section $u_t^\theta=e^{i\theta N_t}$ has norm one and multiplies under independent increments. Products of these units on a partition yield all functions $\exp(i\sum_j\theta_j\Delta_jN)$ of the interval counts.

\paragraph{Totality of the generated vectors.} If $F$ is orthogonal to every such function for a fixed partition, then its conditional expectation given the count vector has zero discrete Fourier transform. The coefficients are absolutely summable by Cauchy--Schwarz against the count probabilities. Multiplying by the opposite character and integrating each angle over $[0,2\pi]$ recovers each coefficient, so the conditional expectation is zero. Rational-time partitions generate the Poisson path sigma-field because the paths are right-continuous step functions. Approximation by functions measurable on successively finer finite partitions then implies $F=0$. Hence products of units are total and the associated system is type I.

\paragraph{What this says about the source.} The ordinary integer-valued finite-rate process cannot furnish the requested type II example. The supplied definition records a published positive answer for certain infinite-dimensional Banach-valued processes; that source qualification is retained, not reclassified as an open broad existence question or claimed as our discovery. The argument above uses recoverability of all finite-partition count functions from these simple units and does not analyze the cited infinite-dimensional construction. The narrower circle-homeomorphism-group example and a reconstruction of the source's type-II mechanism remain unresolved here.

\paragraph{Verification and adversarial audit.}
Checked norm-one multiplicative sections, Fourier coefficient recovery and partition density; preserved the source-qualified already-positive broad existence result. Model self-audit only; no independent human proof checking. Prompts A and B were applied to this canonical record and its qualified source scope. Checked norm-one multiplicative sections, Fourier coefficient recovery and partition density; preserved the source-qualified already-positive broad existence result. Retained unresolved status and zero original-target progress; A type-II construction/reconstruction in the suggested homeomorphism group is not obtained.

Coordinator checked the Poisson character units and totality via discrete Fourier recovery of finite-partition count functions; the supplied broad positive-result qualification remains explicit. This additional check was model-assisted, not independent human verification.

\paragraph{Tools, process and contributions.}
Python 3 was used for file preparation and structural checks; no numerical experiment or theorem prover was used. Contribution prompt A was followed by prompt B. GPT-6 Astra with ultra reasoning drafted the research and performed its stated self-audit; any additional coordinator contribution is identified in the verification and audit. A GPT-6-based Codex coordinator assembled the manuscript and checked repository integration; its exact variant and reasoning configuration were not exposed. No human mathematical contribution, independent human proof check or formal proof certification is supplied. Complete derivations and the consulted source references are included here; private reasoning transcripts are not published.

\paragraph{Final status and first remaining gap.}
\textbf{UNRESOLVED}. A type-II construction/reconstruction in the suggested homeomorphism group is not obtained.

\section{Completion Estimate}
\noindent\textbf{COMPLETION: 0\%}
The calculation is an auxiliary or restricted attempt and establishes no newly verified portion of the full source target. A type-II construction/reconstruction in the suggested homeomorphism group is not obtained.

\subsection{Sources}
\begin{itemize}
\item[\hypertarget{source:statement:1}{S1}] ULAM.AI and the UnsolvedMath Contributors, UnsolvedMath, catalogue entry. Catalogue curation under CC BY 4.0; definitions and English exposition added locally.
\url{https://huggingface.co/datasets/ulamai/UnsolvedMath}.
\item[\hypertarget{source:statement:2}{S2}] Mini-Workshop: Product Systems and Independence in Quantum Dynamics (2009), Oberwolfach Reports, DOI 10.4171/owr/2009/09, pp. 493--548; official report source and mathematical context.
\url{https://doi.org/10.4171/owr/2009/09}.
\item[\hypertarget{source:statement:3}{S3}] Official EMS Press PDF of the same Oberwolfach report; PDF page 53, printed page 545; Floricel--Wadel publisher abstract/introduction and Theorems 5.1--5.2 summaries accessed 2026-10-11.
\url{https://ems.press/content/serial-article-files/46209}.
\item[\hypertarget{source:statement:4}{S4}] Floricel--Wadel 2025; reference in the saved dated catalogue literature review.
\url{https://doi.org/10.1016/j.jfa.2025.111087}.
\item[\hypertarget{source:statement:5}{S5}] R. Floricel and P. Wadel, Product systems arising from Lévy processes, Journal of Functional Analysis 289 (2025), article 111087; primary publisher abstract/introduction consulted.
\url{https://www.sciencedirect.com/science/article/abs/pii/S0022123625002691}.
\end{itemize}
\label{end:statement}
\end{document}
